\documentclass[11pt]{article}
\usepackage{amsmath,amssymb}
\usepackage{geometry}
\geometry{margin=1in}

\title{Equivalence Completion for Longitudinal Gauge Modes and Goldstone Degrees}
\author{Amos Jay Maley}
\date{}

\begin{document}
\maketitle

\section*{Abstract}
We complete an equivalence map between field-level gauge bookkeeping and on-shell high-energy scattering bookkeeping in the electroweak sector. Apparent independent structures—longitudinal vector polarizations and Goldstone scalar modes—are shown to overcount tensor load once admissible redescription and regime transport are enforced. A single invariant longitudinal interaction amplitude class survives. The result is eliminative, invariant under admissible redescription, regime-scoped, and introduces no new physics.

\section{Apparent Independence in Gauge-Field Bookkeeping}
In gauge-fixed electroweak bookkeeping, massive vector fields admit a decomposition into transverse and longitudinal components,
\begin{equation}
W_\mu = W_\mu^{T} + W_\mu^{L},
\end{equation}
with additional scalar fields $\pi^\pm,\pi^0$ appearing as Goldstone modes associated with symmetry breaking. These objects are commonly treated as distinct degrees of freedom at the field level.

\section{Admissible On-Shell Bookkeeping}
An admissible description of scattering processes is provided by on-shell S-matrix elements. In this bookkeeping, invariant content is carried by on-shell amplitudes. Field-level distinctions that do not survive on-shell reduction have no standing as tensor degrees.

\section{Lemma 1: High-Energy Transport}
In the regime $E \gg m_W$, admissible transport from field-level bookkeeping to on-shell scattering yields
\begin{equation}
\mathcal{A}(W_L^\pm, Z_L, \ldots)
\;\longleftrightarrow\;
\mathcal{A}(\pi^\pm, \pi^0, \ldots),
\end{equation}
up to kinematic factors suppressed by $m_W/E$. The transport identifies the same invariant longitudinal interaction content in both descriptions.

\section{Lemma 2: Invariant Witness Counting}
If a tensor load contains $n$ independent degrees of freedom, admissible redescription and regime transport must preserve $n$ independent invariant witnesses across all admissible skins describing the same regime. Distinctions that vanish under admissible transport are redundancy.

\section{Failure of Independence}
Assume longitudinal vector modes and Goldstone scalar modes constitute independent tensor degrees. Under admissible transport to on-shell scattering, only one invariant amplitude class survives. Preservation of independence would require loss of tensor load, hidden unwitnessed structure, or representational privilege. Each constitutes standing failure.

\section{Forced Tensor Relation}
The admissible longitudinal interaction tensor load is uniquely specified by a single amplitude class. In the admissible high-energy regime,
\begin{equation}
\boxed{\;
\mathcal{A}(W_L^\pm, Z_L,\ldots)
\equiv
\mathcal{A}(\pi^\pm, \pi^0,\ldots)
\;}
\end{equation}
The distinction between longitudinal vector polarizations and Goldstone modes is a skin-level decomposition of the same tensor content.

\section{Corollaries}
\begin{itemize}
\item Gauge fixing redistributes longitudinal content without creating new tensor degrees.
\item Higgs exchange contributes to the same invariant longitudinal amplitude class and does not introduce additional longitudinal invariants.
\item Any framework asserting independently tunable longitudinal and Goldstone scattering invariants in the same admissible regime violates standing unless new tensor load is introduced.
\end{itemize}

\section*{Conclusion}
Equivalence completion collapses apparent multiplicity in longitudinal electroweak degrees to a single invariant interaction structure. The result is eliminative, representation-invariant, regime-scoped, and irreversible under admissible redescription.

\end{document}

